\section{Mass generation at one scale implies the cube condition}\label{sec:massgen}

Theorem~\ref{thm:B} reduces the continuum gap to statement (D), a cube
condition uniform along the tuning. This section reduces (D) in turn to a
hypothesis of a different kind, closer to the physics: a uniform decay of
correlations between the gauge-invariant observables of two slabs of a box
separated by a slab of physical width about $\ell/2$, for every boundary
condition. If the maximal correlation is at most $c<\frac12$, then the cube
condition holds with $\varpi=(1-2c)^4/32$ (Theorem~\ref{thm:D}). The value
$\frac12$ is the exact limit of the method (Proposition~\ref{prop:gausschain}
in Appendix~\ref{app:obstructions}); the hypothesis implies exponential decay
of gauge-invariant correlations under every boundary condition
(Proposition~\ref{prop:alldecay}); and it holds in two dimensions along the
weak-coupling tuning (Theorem~\ref{thm:twod}). In four dimensions it is a
form of mass generation at the physical scale and it is not proved here.

\subsection{Boxes, boundary conditions and gauge-invariant observables}

In this section $G\subset\U(N)$ is compact and nontrivial, $\varphi$ is a
positive continuous class function on $G$, $\lambda$ denotes normalised Haar
measure on $G$ or on a product of copies of $G$, and
$\omega_\varphi=\log(\max\varphi/\min\varphi)$. Links of $\Z^4$ are written
$(y,\nu)$ with base point $y$, from $y$ to $y+e_\nu$. A \emph{box} is a set
$Q=\prod_{i=1}^4[\alpha_i,\alpha_i+s_i)\cap\Z$ with integers $s_i\ge1$, its
\emph{widths}, and $V(Q)$ is the set of links with base point in $Q$. For an
exterior configuration $\omega\in G^{E(\Z^4)\setminus V(Q)}$ let
$\nu^\omega_Q$ be the probability measure on $G^{V(Q)}$ with density
proportional to $\prod_{p\cap V(Q)\ne\emptyset}\varphi((\sigma,\omega)_p)$
with respect to $\lambda$, a finite product. It is the conditional law of
$U_{V(Q)}$ given the other links under any Gibbs measure with weight
$\varphi$, and on a torus $\Lambda_L$ with all $s_i\le L-3$ it is the
conditional law of the torus measure.\footnote{The plaquettes meeting $V(Q)$
have their links based in $\prod_i[\alpha_i-1,\alpha_i+s_i]$, which embeds in
$\Lambda_L$ when $s_i\le L-3$.} For $A\subset V(Q)$ write $E_A$ for the
conditional expectation under $\nu^\omega_Q$ given $\sigma_{V(Q)\setminus A}$;
its regular version is integration against the measure on $G^A$ with density
proportional to $\prod_{p\cap A\ne\emptyset}\varphi(U_p)$, which is jointly
continuous in all variables. In particular, if $Q'\subset Q$ is a box, the
conditional law of $\sigma_{V(Q')}$ under $\nu^\omega_Q$ given
$\sigma_{V(Q)\setminus V(Q')}=\eta$ is $\nu^{(\eta,\omega)}_{Q'}$.

A vertex $v$ is \emph{interior} to $Q$ if the eight links incident to $v$ lie
in $V(Q)$. Let $\mathcal G_Q$ be the group of $h\in G^{\Z^4}$ with $h_v=1$ at
every vertex not interior to $Q$, acting by $(h\sigma)_{(y,\nu)}=h_y
\sigma_{(y,\nu)}h_{y+e_\nu}^{-1}$; it fixes every link outside $V(Q)$, whose
end points are not interior. Let $\mathcal A^\omega_Q$ be the set of real
$F\in L^2(\nu^\omega_Q)$ with $F\circ h=F$ almost surely for every $h\in
\mathcal G_Q$: the gauge-invariant observables of the box, the boundary
condition having fixed the gauge on the boundary.

\begin{lemma}\label{lem:boxinv}
Let $Q$ be a box, $\omega$ an exterior configuration and $\nu=\nu^\omega_Q$.
\begin{enumerate}[label=\textup{(\roman*)},leftmargin=*]
\item $\nu$ is invariant under $\mathcal G_Q$.
\item For every $S\subset V(Q)$, each $h\in\mathcal G_Q$ maps
$\sigma(\sigma_S)$ onto itself, and $E_A(F\circ h)=(E_AF)\circ h$ for
$A\subset V(Q)$ and $F\in L^2(\nu)$.
\item $\mathcal A^\omega_Q$ is a closed subspace that contains the constants
and is preserved by every $E_A$. Every $F\in\mathcal A^\omega_Q$ has a
version $\tilde F$ with $\tilde F(h\sigma)=\tilde F(\sigma)$ for all $h$ and
$\sigma$, and if $F$ is $\sigma(\sigma_S)$-measurable, $\tilde F$ can be taken
$\sigma(\sigma_S)$-measurable.
\end{enumerate}
\end{lemma}

\begin{proof}
(i) $h$ multiplies each link on the left and on the right by fixed elements,
so it preserves $\lambda$. Since $h$ fixes the links outside $V(Q)$,
$(h\sigma,\omega)=h\cdot(\sigma,\omega)$, and the holonomy of a plaquette
with base point $y$ is conjugated by $h_y$; $\varphi$ is a class function,
so the density is invariant.

(ii) $(h\sigma)_e$ is a function of $\sigma_e$ alone and $h^{-1}\in\mathcal
G_Q$, so $h$ maps $\sigma(\sigma_S)$ onto itself. For a bounded
$\sigma(\sigma_{V(Q)\setminus A})$-measurable $K$, by (i),
$\nu((F\circ h)K)=\nu(F\,(K\circ h^{-1}))=\nu((E_AF)(K\circ h^{-1}))=
\nu(((E_AF)\circ h)K)$, and $(E_AF)\circ h$ is
$\sigma(\sigma_{V(Q)\setminus A})$-measurable.

(iii) By (i) each map $F\mapsto F\circ h$ is unitary, so its fixed space is
closed, and so is the intersection of these spaces; (ii) gives the
invariance under $E_A$. The group $\mathcal G_Q\cong G^I$, $I$ the finite set
of interior vertices, is compact and acts continuously. For a measurable
version $F_0$ put $\tilde F(\sigma)=\int F_0(h\sigma)\,dh$ where
$\int|F_0(h\sigma)|\,dh<\infty$, and $\tilde F(\sigma)=0$ elsewhere. By the
invariance of Haar measure under right translations, the set where the
integral converges and $\tilde F$ are invariant under every $k\in\mathcal
G_Q$. For each $h$, $F_0(h\sigma)=F_0(\sigma)$ for almost every $\sigma$, so
by Fubini's theorem, for almost every $\sigma$ the equality holds for almost
every $h$, and $\tilde F=F_0$ almost surely. If $F_0$ is a function of
$\sigma_S$, so is $\sigma\mapsto F_0(h\sigma)$ for each $h$, and so is
$\tilde F$.
\end{proof}

\begin{definition}\label{def:rho}
Let $Q$ be a box, $i$ a coordinate, and $Q=Q_-\cup Q_0\cup Q_+$ a splitting
into three nonempty boxes, consecutive in the coordinate $i$ in this order:
the \emph{outer parts} $Q_\pm$ and the \emph{middle} $Q_0$. For an exterior
configuration $\omega$ put $\nu=\nu^\omega_Q$ and
\[
\rho^\omega(Q_-,Q_+)=\sup|\nu(XY)|,
\]
the supremum over $X\in\mathcal A^\omega_Q$ measurable with respect to
$\sigma_{V(Q_-)}$ and $Y\in\mathcal A^\omega_Q$ measurable with respect to
$\sigma_{V(Q_+)}$ with $\nu(X)=\nu(Y)=0$ and $\nu(X^2)=\nu(Y^2)=1$ (the
supremum of the empty set being $0$).
\end{definition}

This is the maximal correlation\footnote{The maximal correlation of two
random variables, the supremum of the correlations of square-integrable
functions of each, was introduced by Hirschfeld and by
Gebelein~\cite{Geb41}; for jointly Gaussian variables it equals the
absolute value of the ordinary correlation.} between the gauge-invariant
observables of the two outer parts, the middle being integrated out, under
the boundary condition $\omega$. By the Cauchy--Schwarz inequality
$\rho^\omega\le1$.

\begin{definition}\label{def:M}
Let $d\ge2$, $m\ge4$ and $w\ge1$ be integers and $0\le c\le1$.
\begin{enumerate}[label=\textup{(\roman*)},leftmargin=*]
\item $(\mathrm M_{d,m}(c))$: $\rho^\omega(Q_-,Q_+)\le c$ for every box $Q$
all of whose widths lie in $[2d-1,md-1]$, every coordinate $i$ in which $Q$
has width $3d-1$, the splitting in which $Q_-$, $Q_0$, $Q_+$ have widths
$d$, $d-1$, $d$ in the coordinate $i$, and every exterior configuration
$\omega$.
\item $(\mathrm M^*_w(c))$: $\rho^\omega(Q_-,Q_+)\le c$ for every box $Q$,
every coordinate $i$, every splitting whose middle has width $w$ in the
coordinate $i$, the outer parts being of any widths, and every $\omega$.
\end{enumerate}
\end{definition}

The boxes of $(\mathrm M_{d,m}(c))$ are fat: all their widths are at least
$2d-1$. The restriction matters. In two dimensions the same condition for
boxes of width one in a transverse direction fails along the weak-coupling
tuning, through a broken global symmetry of a pinned strip, while the fat
form holds there (Theorem~\ref{thm:twod} and Proposition~\ref{prop:strip}
in Appendix~\ref{app:obstructions}). At $\varphi\equiv1$ the links are
independent under $\nu^\omega_Q$, so $\rho^\omega=0$ and both conditions hold
with $c=0$. Along the tuning, with $d\approx\ell/(2a)$, $(\mathrm
M_{d,m}(c))$ is a statement at the physical separation $(d-1)a\approx\ell/2$
in boxes of physical size at most about $m\ell/2$, uniform in the boundary
condition, and stated for all square-integrable gauge-invariant observables
of the two regions rather than for pairs of local ones.

\subsection{Two projections and a chain of projections}

\begin{lemma}\label{lem:friedrichs}
Let $P_1,P_2$ be orthogonal projections on a real Hilbert space $\mathcal K$,
$\mathcal N$ a closed subspace of $\operatorname{ran}P_1\cap
\operatorname{ran}P_2$, and suppose that $|\ip xy|\le c\|x\|\|y\|$ for all
$x\in\operatorname{ran}P_1\ominus\mathcal N$ and $y\in\operatorname{ran}P_2
\ominus\mathcal N$. Put $Q_k=I-P_k$ and $X=Q_1+Q_2$. Then $X^2\ge(1-c)X$,
that is, $Q_1Q_2+Q_2Q_1\ge-c(Q_1+Q_2)$.
\end{lemma}

\begin{proof}
For $n\in\mathcal N$, $P_kn=n$, so $Xn=0$; as $X$ is self-adjoint it also
preserves $\mathcal N^\perp$. Let $g\in\mathcal N^\perp$. Then $P_kg\perp
\mathcal N$, since $\ip{P_kg}n=\ip g{P_kn}=\ip gn=0$, so with $p=\|P_1g\|$
and $q=\|P_2g\|$, $|\ip{P_1g}{P_2g}|\le cpq$. Put $s=\ip g{(P_1+P_2)g}=
p^2+q^2$. Then $\|(P_1+P_2)g\|^2=p^2+q^2+2\ip{P_1g}{P_2g}\le(1+c)(p^2+q^2)$
and $s\le\|g\|\|(P_1+P_2)g\|\le\|g\|((1+c)s)^{1/2}$, so $s\le(1+c)\|g\|^2$
and $\ip g{Xg}=2\|g\|^2-s\ge(1-c)\|g\|^2$. So $X=0$ on $\mathcal N$, and on
$\mathcal N^\perp$ the operator $X-(1-c)$ is nonnegative and commutes with
$X\ge0$; hence $X^2-(1-c)X=X^{1/2}(X-(1-c))X^{1/2}\ge0$ there. Finally
$X^2-X=Q_1Q_2+Q_2Q_1$.
\end{proof}

The angle between $\operatorname{ran}P_1\ominus\mathcal N$ and
$\operatorname{ran}P_2\ominus\mathcal N$ is the Friedrichs angle of the two
ranges, and the lemma is a form of the classical relation between that angle
and the rate of the alternating projections~\cite{Deu01}.

\begin{lemma}\label{lem:chain}
Let $Q_0,\dots,Q_M$ be orthogonal projections on a real Hilbert space
$\mathcal K$ with $Q_tQ_{t'}=Q_{t'}Q_t$ whenever $|t-t'|\ge2$ and
$Q_tQ_{t+1}+Q_{t+1}Q_t\ge-c(Q_t+Q_{t+1})$ for $0\le t<M$, with $0\le c<
\frac12$. Then $H=\sum_tQ_t$ satisfies $H^2\ge(1-2c)H$, and $\ip g{Hg}\ge
(1-2c)\|g\|^2$ for every $g\perp\ker H$.
\end{lemma}

\begin{proof}
$H^2=\sum_tQ_t+\sum_{t<M}(Q_tQ_{t+1}+Q_{t+1}Q_t)+\sum Q_tQ_{t'}$, the last
sum over the ordered pairs with $|t-t'|\ge2$. Each such $Q_tQ_{t'}$ is a
product of commuting projections, hence a projection and nonnegative. The
middle sum is at least $-c\sum_{t<M}(Q_t+Q_{t+1})\ge-2cH$. So $H^2-(1-2c)H
\ge0$, and by the spectral theorem the spectrum of the nonnegative operator
$H$ lies in $\{0\}\cup[1-2c,\infty)$, which gives the second claim.
\end{proof}

\begin{lemma}\label{lem:constant}
Let $Q$ be a box, $\omega$ an exterior configuration, and $R_0,\dots,R_M$
subsets of $V(Q)$ with union $V(Q)$. If $F\in L^2(\nu^\omega_Q)$ is, for every
$t$, almost surely equal to a $\sigma(\sigma_{V(Q)\setminus
R_t})$-measurable function, then $F$ is almost surely constant.
\end{lemma}

\begin{proof}
The density of $\nu=\nu^\omega_Q$ with respect to $\lambda$ is bounded above
and below by positive constants, so $L^2(\nu)=L^2(\lambda)$ with the same
null sets. Let $\Lambda_t$ be the conditional expectation under $\lambda$
given $\sigma_{V(Q)\setminus R_t}$, which integrates out $\sigma_{R_t}$. By
hypothesis $F=F_t$ almost surely with $F_t$ a function of
$\sigma_{V(Q)\setminus R_t}$, so $\Lambda_tF=F_t=F$. By Fubini's theorem
$\Lambda_0\cdots\Lambda_M$ integrates out all of $\bigcup_tR_t=V(Q)$, so
$F=\Lambda_0\cdots\Lambda_MF$ is constant.
\end{proof}

\subsection{One direction}

Fix a box $R=\prod_iI_i$ and a coordinate $i$ with $I_i=[\alpha,\alpha+
(M+2)d-1)$ for an integer $M\ge0$. For $0\le t\le M$ let $R_t$ be the box
obtained from $R$ by replacing $I_i$ with $\alpha+dt+[0,2d-1)$, and for
$t<M$ let $Q_t$ be the box obtained by replacing $I_i$ with $\alpha+dt+
[0,3d-1)$, split into $Q_t^-$, $Q_t^0$, $Q_t^+$ by $\alpha+dt+[0,d)$,
$\alpha+dt+[d,2d-1)$ and $\alpha+dt+[2d-1,3d-1)$. Then
\begin{equation}\label{eq:slabs}
V(Q_t)=V(R_t)\cup V(R_{t+1}),\quad V(Q_t^-)=V(R_t)\setminus V(R_{t+1}),
\quad V(Q_t^+)=V(R_{t+1})\setminus V(R_t),
\end{equation}
and $\bigcup_tV(R_t)=V(R)$, since the intervals $\alpha+dt+[0,2d-1)$ cover
$I_i$.

\begin{proposition}[one direction]\label{prop:slab}
In this setting let $\omega$ be an exterior configuration of $R$,
$\nu=\nu^\omega_R$, $P_t=E_{V(R_t)}$, and suppose that $\rho^{\omega'}(Q^-_t,
Q^+_t)\le c<\frac12$ for every $t<M$ and every exterior configuration
$\omega'$ of $Q_t$. Then every $F\in\mathcal A^\omega_R$ satisfies
\[
(1-2c)\Var_\nu(F)\le\sum_{t=0}^M\nu\bigl((F-P_tF)^2\bigr).
\]
\end{proposition}

\begin{proof}
For $M=0$, $P_0F=\nu(F)$ and the claim is trivial; let $M\ge1$. Put
$\mathcal K=\mathcal A^\omega_R$. By Lemma~\ref{lem:boxinv}(iii) each $P_t$
preserves $\mathcal K$ and restricts to an orthogonal projection of
$\mathcal K$; let $Q_t=I-P_t$ on $\mathcal K$.

\emph{Non-adjacent moves commute.} Let $t'\ge t+2$. The base points of the
links of $V(R_t)$ have $i$-th coordinate at most $\alpha+dt+2d-2$ and those
of $V(R_{t'})$ at least $\alpha+dt+2d$, so no plaquette meets both sets.
Given $\sigma$ off $V(R_t)\cup V(R_{t'})$, the conditional density of
$(\sigma_{V(R_t)},\sigma_{V(R_{t'})})$ is a product of factors each of which
involves the links of only one of the two sets. So they are conditionally
independent, and $P_tP_{t'}F$ and $P_{t'}P_tF$ are both the integral of $F$
against the product of the two conditional laws. Hence $P_tP_{t'}=
P_{t'}P_t$, and the $Q_t$ commute.

\emph{Adjacent moves.} Fix $t<M$ and write $a=V(Q_t^-)$, $b=V(Q_t^0)$,
$c'=V(Q_t^+)$, $e=V(R)\setminus V(Q_t)$. By~\eqref{eq:slabs},
$V(R)\setminus V(R_t)=c'\cup e$ and $V(R)\setminus V(R_{t+1})=a\cup e$, so
the ranges of $P_t$ and of $P_{t+1}$ in $\mathcal K$ are the elements of
$\mathcal K$ measurable with respect to $(\sigma_{c'},\sigma_e)$ and to
$(\sigma_a,\sigma_e)$. Let $\mathcal N$ be the elements of $\mathcal K$
measurable with respect to $\sigma_e$, a closed subspace of both ranges. We
check the hypothesis of Lemma~\ref{lem:friedrichs}. Let $x$ and $y$ be in
the ranges of $P_t$ and $P_{t+1}$ and orthogonal to $\mathcal N$. The
conditional expectation $E_e$ given $\sigma_e$ is $E_{V(Q_t)}$, which
preserves $\mathcal K$, so $E_ex\in\mathcal N$ and $\|E_ex\|^2=\ip x{E_ex}=0$;
likewise $E_ey=0$. By Lemma~\ref{lem:boxinv}(iii) for the box $R$ take
versions $x=\tilde x(\sigma_{c'},\sigma_e)$ and $y=\tilde y(\sigma_a,\sigma_e)$
invariant at every point under $\mathcal G_R$. Every vertex interior to $Q_t$
is interior to $R$, so $\mathcal G_{Q_t}\subset\mathcal G_R$, and
$\mathcal G_{Q_t}$ fixes $\sigma_e$. Hence for each $\eta\in G^e$ the
functions $X_\eta=\tilde x(\cdot,\eta)$ of $\sigma_{V(Q^+_t)}$ and $Y_\eta=
\tilde y(\cdot,\eta)$ of $\sigma_{V(Q^-_t)}$ are invariant under
$\mathcal G_{Q_t}$ at every point. The conditional law of $\sigma_{V(Q_t)}$
given $\sigma_e=\eta$ is $\nu_\eta=\nu^{(\eta,\omega)}_{Q_t}$, so for almost
every $\eta$, $X_\eta,Y_\eta\in L^2(\nu_\eta)$, $\nu_\eta(X_\eta)=(E_ex)(\eta)
=0$, $\nu_\eta(Y_\eta)=0$, and by the hypothesis with $\omega'=(\eta,\omega)$,
\[
|\nu_\eta(X_\eta Y_\eta)|\le c\,\nu_\eta(X_\eta^2)^{1/2}\nu_\eta(Y_\eta^2)^{1/2}.
\]
Integrating over $\eta$ and using the Cauchy--Schwarz inequality,
$|\ip xy|\le c\|x\|\|y\|$, and Lemma~\ref{lem:friedrichs} gives
$Q_tQ_{t+1}+Q_{t+1}Q_t\ge-c(Q_t+Q_{t+1})$ on $\mathcal K$.

\emph{Conclusion.} Lemma~\ref{lem:chain} applies. If $Hg=0$ for $g\in
\mathcal K$, then $\sum_t\|Q_tg\|^2=\ip g{Hg}=0$, so $g=P_tg$ for all $t$ and
$g$ is constant by Lemma~\ref{lem:constant}. For $F\in\mathcal K$ the
function $g=F-\nu(F)$ is orthogonal to the constants, hence to $\ker H$, and
$\ip g{Hg}=\sum_t\nu((F-P_tF)^2)$.
\end{proof}

\subsection{The patch gap and the cube condition}

For $d\ge2$, $m\ge4$ and $r\in\Z^4$ let $\Pi=r+[0,md-1)^4$, the box of
base points of a patch, and for $j\in[0,m-2]^4$ let $B_j=V(r+dj+
[0,2d-1)^4)$, as in Appendix~\ref{app:finite}. For an exterior
configuration $\omega$ of $\Pi$ write $\nu=\nu^\omega_\Pi$ and $P_j=E_{B_j}$.

\begin{theorem}[patch gap from slab decorrelation]\label{thm:patch}
Let $d\ge2$, $m\ge4$, $0\le c<\frac12$, and assume $(\mathrm M_{d,m}(c))$.
Then for every $r$, every $\omega$ and every $F\in\mathcal A^\omega_\Pi$,
\[
(1-2c)^4\Var_\nu(F)\le\sum_{j\in[0,m-2]^4}\nu\bigl((F-P_jF)^2\bigr).
\]
\end{theorem}

\begin{proof}
For $0\le k\le4$ and $\tau\in[0,m-2]^k$ let $R^{(k)}_\tau=r+\prod_{i\le k}
(d\tau_i+[0,2d-1))\times[0,md-1)^{4-k}$ and $P^{(k)}_\tau=
E_{V(R^{(k)}_\tau)}$, under $\nu$. We prove by induction on $k$ that
\[
(1-2c)^k\Var_\nu(F)\le\sum_{\tau\in[0,m-2]^k}\nu\bigl((F-P^{(k)}_\tau F)^2
\bigr)\qquad(F\in\mathcal A^\omega_\Pi).
\]
For $k=0$, $R^{(0)}=\Pi$ and $P^{(0)}F=\nu(F)$. Let $k<4$, fix $\tau$, write
$R=R^{(k)}_\tau$, and split $\sigma=(\sigma',\omega')$ with $\sigma'=
\sigma_{V(R)}$ and $\omega'=\sigma_{V(\Pi)\setminus V(R)}$. Take $F$
invariant at every point under $\mathcal G_\Pi$ (Lemma~\ref{lem:boxinv}(iii)).
Since $\mathcal G_R\subset\mathcal G_\Pi$ fixes $\omega'$, each
$F(\cdot,\omega')$ is invariant under $\mathcal G_R$. The conditional law of
$\sigma'$ given $\omega'$ is $\nu'=\nu^{(\omega',\omega)}_R$, so
\[
\nu\bigl((F-P^{(k)}_\tau F)^2\bigr)=\int\Var_{\nu'}\bigl(F(\cdot,\omega')
\bigr)\,d\nu(\omega').
\]
In the coordinate $k+1$ the box $R$ has width $md-1=((m-2)+2)d-1$. The
boxes $R_u$ and $Q_u$ of Proposition~\ref{prop:slab} built from $R$ with
$i=k+1$ and $M=m-2$ are $R_u=R^{(k+1)}_{(\tau,u)}$ and boxes whose widths are
$2d-1$ in the coordinates $\le k$, $3d-1$ in the coordinate $k+1$ and $md-1$
in the coordinates $>k+1$, split as in $(\mathrm M_{d,m}(c))$; all these
widths lie in $[2d-1,md-1]$. By Proposition~\ref{prop:slab}, for almost every
$\omega'$,
\[
(1-2c)\Var_{\nu'}(F(\cdot,\omega'))\le\sum_{u=0}^{m-2}\nu'\bigl((F(\cdot,
\omega')-P'_uF(\cdot,\omega'))^2\bigr),
\]
with $P'_u$ the conditional expectation under $\nu'$ given $\sigma'$ off
$V(R^{(k+1)}_{(\tau,u)})$. The conditional law of the links of that set given
the other links of $\Pi$ is the same under $\nu$ and under $\nu'$, both
having density proportional to $\prod_p\varphi(U_p)$ over the plaquettes
meeting it; so $P'_uF(\cdot,\omega')=(P^{(k+1)}_{(\tau,u)}F)(\cdot,\omega')$,
and integrating over $\omega'$ gives $\nu((F-P^{(k+1)}_{(\tau,u)}F)^2)$.
Multiplying the induction hypothesis by $1-2c$ and inserting these bounds
proves the case $k+1$. For $k=4$, $V(R^{(4)}_\tau)=B_\tau$.
\end{proof}

\begin{mainthm}[mass generation implies the cube condition]\label{thm:D}
Let $d\ge2$, $0\le c<\frac12$ and $m\ge4$ with $m-1\ge148/(1-2c)^4$, and
assume $(\mathrm M_{d,m}(c))$. Then for every even $L$ with $d\mid L$ and
$L\ge2md$ the measure $\mu_L$ with weight $\varphi$ satisfies
\[
\BP\Bigl(2d-1,\frac{(1-2c)^4}{32}\Bigr).
\]
Consequently, let $(\beta_k,a_k)$ be a tuning with $a_k\to0$ and Wilson
weights, and $d_k\ge2$ integers with $(2d_k-1)a_k\le\ell$. If
$(\mathrm M_{d_k,m}(c))$ holds for every $k$ with $c<\frac12$ and $m$ as
above independent of $k$, then statement (D) holds with $n_k=2d_k-1$ and
$\varpi=(1-2c)^4/32$, and every continuum limit as in Theorem~\ref{thm:B}
along the tori with $d_k\mid L$ satisfies
\[
\sigma(H)\subset\{0\}\cup\Bigl[\frac{(1-2c)^4}{1250\,e^2\ell},\infty\Bigr),
\]
with $0$ a simple eigenvalue; a nontrivial limit has a mass gap at least
this large.
\end{mainthm}

\begin{proof}
Theorem~\ref{thm:patch} says that the gauge-invariant patch gap of
Definition~\ref{def:patchgap} satisfies $\gamma_{d,m}\ge(1-2c)^4$, and
Theorem~\ref{thm:finite} with $\gamma=(1-2c)^4$ gives
$\BP(2d-1,(1-2c)^4/32)$. Along the tuning this holds on an infinite set of
even tori with $L\ge2(2d_k-1)+4$ for each $k$, with $(2d_k-1)a_k\le\ell$, and
Theorem~\ref{thm:B} gives the gap $16\varpi/(625e^2\ell)$ with
$\varpi=(1-2c)^4/32$.
\end{proof}

The patch gap obtained, $(1-2c)^4$, does not decay with $m$, whereas the
finite-size criterion needs only $148/(m-1)$; this slack is what allows a
large but fixed number $m$ of blocks per patch. The threshold $c<\frac12$
cannot be raised for this argument: for the massless Gaussian chain every
adjacent pair of moves has maximal correlation exactly $\frac12$, the
non-adjacent moves commute, and the gap of the slab dynamics tends to zero
(Proposition~\ref{prop:gausschain}).

\subsection{Decay of correlations under every boundary condition}

\begin{lemma}\label{lem:boxmarkov}
Let $Q$ be a box, $\omega$ an exterior configuration, $i$ a coordinate, and
$Y_0,\dots,Y_K$ the sets $V(Q\cap\{x_i\in J_k\})$ for consecutive nonempty
integer intervals $J_0<J_1<\dots<J_K$ whose union is the $i$-th side of $Q$.
Under $\nu^\omega_Q$, for each $k$, the families $(\sigma_{Y_l})_{l<k}$ and
$(\sigma_{Y_l})_{l>k}$ are conditionally independent given $\sigma_{Y_k}$.
\end{lemma}

\begin{proof}
The links in $V(Q)$ of a plaquette are based at points whose $i$-th
coordinates take at most two consecutive values, so they lie in one $Y_l$ or
in two consecutive ones. Grouping the factors of the density by the least
$l$ met, the density of $\nu^\omega_Q$ with respect to $\lambda$ is
$\prod_l\psi_l(\sigma_{Y_l},\sigma_{Y_{l+1}})$, with $\psi_K$ a function of
$\sigma_{Y_K}$, which gives the conditional independence.
\end{proof}

\begin{proposition}\label{prop:alldecay}
Assume $(\mathrm M^*_w(c))$ with $w\ge1$ and $0\le c\le1$. Let $Q$ be a box,
$\omega$ an exterior configuration, $\nu=\nu^\omega_Q$, $i$ a coordinate,
$s\in\Z$, $r\ge1$, and let $f,g\in\mathcal A^\omega_Q$ be measurable with
respect to $\sigma_{V(Q\cap\{x_i\le s\})}$ and $\sigma_{V(Q\cap\{x_i\ge
s+r\})}$. Then
\[
|\Cov_\nu(f,g)|\le c^{\lfloor r/(w+1)\rfloor}\|f-\nu(f)\|\,\|g-\nu(g)\|.
\]
\end{proposition}

\begin{proof}
If one of the two parts of $Q$ is empty, $f$ or $g$ is constant and the
covariance vanishes. Otherwise the $i$-th side of $Q$ contains $[s,s+r]$.
Let $J=\lfloor r/(w+1)\rfloor$; for $J=0$ the claim is the Cauchy--Schwarz
inequality. Let $J\ge1$ and split the $i$-th side into consecutive
intervals: $T_0$, the values $\le s$; for $1\le k\le J$, $M_k=[s+(k-1)(w+1)+
1,\,s+(k-1)(w+1)+w]$; for $1\le k\le J-1$, $T_k=\{s+k(w+1)\}$; and $T_J$, the
values $\ge s+J(w+1)$, which contain the values $\ge s+r$. Use the same
letters for the corresponding sets of links. By Lemma~\ref{lem:boxmarkov}
the sequence $T_0,M_1,T_1,\dots,M_J,T_J$ is Markov, and by
Lemma~\ref{lem:boxinv} the conditional expectations given $\sigma_{T_k}$
preserve $\mathcal A^\omega_Q$.

For $1\le k\le J$ let $u\in\mathcal A^\omega_Q$ be a centred function of
$\sigma_{T_{k-1}}$ and $v=E[u\mid\sigma_{T_k}]$, which is centred and lies
in $\mathcal A^\omega_Q$. The box $Q$ splits in the coordinate $i$ into the
values $\le s+(k-1)(w+1)$, the middle $M_k$ of width $w$, and the values
$\ge s+k(w+1)$; $u$ is a function of the first part and $v$ of the last, so
by $(\mathrm M^*_w(c))$, $\|v\|^2=\nu(uv)\le c\|u\|\|v\|$, that is,
$\|v\|\le c\|u\|$. Let $f_0=f-\nu(f)$ and $u_k=E[f_0\mid\sigma_{T_k}]$. By
the Markov property $u_k=E[u_{k-1}\mid\sigma_{T_k}]$, so $\|u_k\|\le
c\|u_{k-1}\|$ and $\|u_J\|\le c^J\|f_0\|$. As $g$ is
$\sigma_{T_J}$-measurable, $\Cov_\nu(f,g)=\nu(u_J(g-\nu(g)))$, and the
Cauchy--Schwarz inequality gives the claim.
\end{proof}

So $(\mathrm M^*_{d-1}(c))$ with $c<1$, which contains $(\mathrm
M_{d,m}(c))$ for every $m$, gives exponential decay of the covariances of
gauge-invariant observables of a box, uniformly in the box and in the
boundary condition, at rate $\log(1/c)/d$ per lattice spacing, that is, of
order $\log(1/c)/\ell$ in physical units along the tuning. This is close to
the decay hypothesis of Chatterjee~\cite[Def.~2.3]{Cha21}, under which an
area law for Wilson loops follows; there the observables are local
functions of a cube with prescribed boundary, and the geometries differ, so
neither condition is claimed to imply the other. The converse direction is
the weak point of any pairwise form: a bound $K_1e^{-K_2r}$ for pairs of
local functions, summed over the pairs of two slabs of transverse size
$(md)^3$, bounds $\rho^\omega$ only up to a factor of order $(md)^6$, which
is not uniform in the lattice spacing.

\subsection{What is known about the hypothesis}

\emph{Strong coupling at fixed $d$ and $m$.} The hypothesis is not
vacuous: $(\mathrm M_{d,m}(c))$ holds with $c=1-e^{-24\omega_\varphi
(md)^4}$. Let $Q$, the splitting and $\omega$ be as in Definition~\ref{def:M},
$A_1=V(Q_-\cup Q_0)$, $A_2=V(Q_0\cup Q_+)$, $\nu=\nu^\omega_Q$, and let
$\gamma_2$ be the best constant in
\[
\gamma_2\Var_\nu(g)\le\sum_{k=1,2}\nu\bigl((g-E_{A_k}g)^2\bigr),\qquad g\in
L^2(\nu)\text{ real}.
\] Under $\lambda$ the links are independent
and $A_1\cup A_2=V(Q)$, so the Efron--Stein inequality~\cite{ES81} gives the
constant $1$ for $\lambda$. The logarithm of the density of $\nu$ with
respect to $\lambda$ has oscillation $\Omega$ at most $\omega_\varphi$ times
the number of plaquettes meeting $V(Q)$, which is at most $24(md)^4$, and a
change of the reference measure by a density whose logarithm has oscillation
$\Omega$ changes the constant by at most the factor
$e^{-\Omega}$~\cite{HS87}.\footnote{Write $d\nu/d\lambda=e^V/Z$. Then $\Var_\nu(g)\le
Z^{-1}e^{\sup V}\Var_\lambda(g)$, and $\lambda((g-E^\lambda_Ag)^2)\le
Ze^{-\inf V}\nu((g-E_Ag)^2)$, as both conditional expectations project onto
the functions of the links outside $A$.} So $\gamma_2\ge e^{-24\omega_\varphi
(md)^4}$. Now let $X,Y$ be as in Definition~\ref{def:rho} with
$r=\nu(XY)\ge0$, and $g=X+Y$, so $\nu(g)=0$ and $\|g\|^2=2+2r$. Since $X$ is a
function of links off $A_2$, $E_{A_2}X=X$, and $\|X-E_{A_1}X\|^2=1-
\|E_{A_1}X\|^2\le1-r^2$ because $Y$ is a unit vector in the range of $E_{A_1}$
and $|\ip{E_{A_1}X}Y|=|\ip XY|=r$; symmetrically for $Y$. Hence
\[
\gamma_2(2+2r)\le\sum_k\|g-E_{A_k}g\|^2\le2(1-r^2),
\]
so $r\le1-\gamma_2$. The range of
couplings obtained in this way shrinks with $md$, and it is far smaller than
that of Theorem~\ref{thm:C}.

\emph{Two dimensions.} In two dimensions the seam between the outer parts
of a box is one-dimensional, and after an axial gauge the middle slab
integrates out exactly to a convolution power of the weight. This proves the
two-dimensional form of the hypothesis along the weak-coupling tuning, with
$c\to0$ as the physical size of the box grows, uniformly in the boundary
condition (Theorem~\ref{thm:twod} in Appendix~\ref{app:obstructions}). In
four dimensions the seam carries a three-dimensional layer of links, and
$(\mathrm M_{d,m}(c))$ is a bound, uniform in the boundary condition, on the
second singular value of the $(d-1)$-step transfer operator of the box on
gauge-invariant functions. Boundary conditions with values in the centre,
't~Hooft's twisted ones~\cite{tH79}, leave a global conjugation symmetry of
the box, and in the two-dimensional strip such a symmetry is broken at weak
coupling (Proposition~\ref{prop:strip}); in four dimensions the hypothesis
asks that the twisted sectors be disordered at the scale $\ell/2$, which is
the form that confinement takes here. We know no boundary condition that
refutes $(\mathrm M_{d,m}(c))$ along the four-dimensional tuning, and no
proof of it.

\emph{An area law does not suffice.} The next proposition shows that the
Wilson loops of the defining representation carry no information about the
observables that (D) constrains.

\begin{proposition}[central extensions]\label{prop:cover}
Let $\tilde G\subset\U(\tilde N)$ and $G\subset\U(N)$ be compact,
$\pi:\tilde G\to G$ a continuous surjective homomorphism with finite central
kernel $Z$, $\varphi$ a positive continuous class function on $G$ and
$\tilde\varphi=\varphi\circ\pi$. On $\Lambda_L$ let $\tilde\mu$, $\mu$ be the
measures with weights $\tilde\varphi$, $\varphi$, and $\Pi:\tilde G^E\to G^E$
the map $\pi$ applied to every link.
\begin{enumerate}[label=\textup{(\roman*)},leftmargin=*]
\item $\tilde\varphi$ is a positive continuous class function and
$\Pi_*\tilde\mu=\mu$.
\item For every $x$ and $F\in L^2(\mu)$, $\tilde{\mathcal P}_x(F\circ\Pi)=
(\mathcal P_xF)\circ\Pi$ and $\Var_{\tilde\mu}(F\circ\Pi)=\Var_\mu(F)$; if $F$
is gauge invariant, so is $F\circ\Pi$. Hence $\BP(n,\varpi)$ for $\tilde\mu$
implies $\BP(n,\varpi)$ for $\mu$.
\item If some $z\in Z$ acts on $\C^{\tilde N}$ as a scalar $\zeta\ne1$, then
$\tilde\mu(\Tr\tilde U_C)=0$ for every closed lattice path $C$ that traverses
some link exactly once.
\end{enumerate}
\end{proposition}

\begin{proof}
(i) $\pi(hgh^{-1})=\pi(h)\pi(g)\pi(h)^{-1}$, so $\tilde\varphi$ is a class
function. The image under $\pi$ of Haar measure on $\tilde G$ is a
left-invariant probability measure on $G$, hence Haar measure, so
$\Pi_*\tilde\lambda=\lambda$. As $\pi$ is a homomorphism, $\pi(\tilde U_p)=
(\Pi\tilde U)_p$, so $\tilde\mu$ has density $\Psi\circ\Pi$ with respect to
$\tilde\lambda$, with $\Psi(U)=\tilde Z^{-1}\prod_p\varphi(U_p)$, and
$\Pi_*((\Psi\circ\Pi)\tilde\lambda)=\Psi\lambda$; both being probability
measures, $\Psi\lambda=\mu$.

(ii) $\tilde{\mathcal P}_x(F\circ\Pi)(\tilde U)$ is the integral of
$F(\Pi\tilde\sigma,\Pi\tilde U_{\rm out})$ against the density proportional to
$\prod_{p\cap V_x\ne\emptyset}\varphi((\Pi\tilde\sigma,\Pi\tilde U_{\rm
out})_p)$ on $\tilde G^{V_x}$; by the change of variables $\sigma=\Pi\tilde
\sigma$ and (i) on $V_x$ it equals $(\mathcal P_xF)(\Pi\tilde U)$. The variance
identity follows from (i), and $\Pi(\tilde g\tilde U)=\pi(\tilde g)\Pi(\tilde
U)$ for a gauge transformation $\tilde g$ gives the invariance. Then for
gauge-invariant $F$, $\varpi\Var_\mu(F)=\varpi\Var_{\tilde\mu}(F\circ\Pi)\le
n^{-4}\sum_x\tilde\mu((F\circ\Pi-\tilde{\mathcal P}_x(F\circ\Pi))^2)=n^{-4}
\sum_x\mu((F-\mathcal P_xF)^2)$.

(iii) Let $e$ be a link traversed once by $C$ and $T$ the map $\tilde U_e
\mapsto z\tilde U_e$, the other links being fixed. It preserves $\tilde
\lambda$, and as $z$ is central with $\pi(z)=1$, each plaquette variable
through $e$ becomes $z^{\pm1}\tilde U_p$ with $\tilde\varphi(z^{\pm1}\tilde
U_p)=\tilde\varphi(\tilde U_p)$; so $\tilde\mu$ is $T$-invariant. Moving the
central factor out of the product, $\Tr\tilde U_C\circ T=\zeta^{\pm1}\Tr
\tilde U_C$, so $\tilde\mu(\Tr\tilde U_C)=\zeta^{\pm1}\tilde\mu(\Tr\tilde U_C)
=0$.
\end{proof}

For instance, with $\tilde G=G=\U(1)$, $\pi(z)=z^2$ and the Wilson weight
$\varphi$ of $\U(1)$, every Wilson loop of charge one under $\tilde\varphi(z)
=e^{\beta\operatorname{Re}z^2}$ that traverses some link once vanishes for
every $\beta$, so Wilson's area law holds with every string tension, while
by part (ii) the cube condition for $\tilde\varphi$ is at least as hard as for
the $\U(1)$ theory, in which it fails at weak coupling in the Coulomb phase
(Section~\ref{sec:weak}). The vanishing is kinematic, a consequence
of a local symmetry and not of confinement; what the example shows is that
an argument deriving (D) from an area law of the defining representation
cannot be valid for all compact groups and weights. For $\SU(N)$ with the
Wilson action no such example is known, and nothing is claimed about whether
an area law and (D) hold together there.
